%%=============================================================================
%% Voorwoord
%%=============================================================================

\chapter*{\IfLanguageName{dutch}{Woord vooraf}{Preface}}%
\label{ch:voorwoord}

%% TODO:
%% Het voorwoord is het enige deel van de bachelorproef waar je vanuit je
%% eigen standpunt (``ik-vorm'') mag schrijven. Je kan hier bv. motiveren
%% waarom jij het onderwerp wil bespreken.
%% Vergeet ook niet te bedanken wie je geholpen/gesteund/... heeft

Als student van Hogent vormt deze bachelorproef de afronding van mijn opleiding toegepaste informatica met afstudeerrichting AI and data engineering. Tijdens mijn opleiding groeide mijn interesse in het vakgebied artificiële intelligentie en mijn leergierigheid om data om te zetten naar waardevolle inzichten. Het idee dat data kan helpen om betere beslissingen te nemen, sprak mij enorm aan. Hierdoor koos ik ervoor om een bachelorproef uit te werken die zowel het domein van artificiële intelligentie verkent als waardevolle inzichten levert aan bedrijven.

Het uitvoeren van deze bachelorproef was een leerrijk maar uitdagend traject. Ik kreeg de kans om mijn theoretische kennis toe te passen in een realistische context en maakte kennis met onderwerpen zoals dataverzameling, model optimalisatie en experimentatie. Daarbij leerde ik niet alleen technisch bij, maar ontwikkelde ik ook vaardigheden zoals zelfstandig werken, plannen en kritisch denken.

Ik wil mijn promotors Koen Mertens en Lena De Mol bedanken voor hun waardevolle feedback en kritische inzichten die mij tijdens dit onderzoek in de juiste richting stuurden. Verder wil ik ook mijn co-promotor Marten Roels bedanken voor zijn steun, constructieve feedback en technische ondersteuning tijdens mijn bachelorproef. Tot slot wil ik in het bijzonder het bedrijf Accurat bedanken voor de steun en middelen die ze ter beschikking stelden doorheen deze bachelorproef, waardoor ik zowel mijn bachelorproef als de proof of concept tot een goed einde kon brengen.

Tot slot hoop ik dat deze bachelorproef een waardevolle bijdrage levert en dat de opgedane kennis een stevige basis vormt voor mijn verdere carrière in AI en data engineering.
